\documentclass[11pt]{article}
\usepackage[margin=1in]{geometry}
\usepackage{amsmath,amssymb,amsthm,hyperref}
\newtheorem{theorem}{Theorem}
\newcommand{\G}{\mathcal G}
\newcommand{\E}{\mathcal E}
\newcommand{\C}{\mathcal C}
\title{Strong triangularity of the Hermitian Jacobi operator\\\large MathDB 376317: a Casimir proof}
\author{Principia Math In Progress}
\date{17 September 2026}
\begin{document}
\maketitle

\paragraph{Task 1: the original problem.}
We prove Kuijper's Conjecture 3.13 \cite{K}, with strong triangularity
understood exactly as in Definition 2.12 and with the ordering on
pp.~19--20. The diagonal coefficient is independent of every choice
within a multiplicity space.

\begin{theorem}\label{thm}
Let $m,k\ge1$ and $\kappa_j>m/2-1$ for $1\le j\le k+1$. Put
\[
 N=\sum_{j=1}^{k+1}\kappa_j+\frac{(k+1)m}{2}.
\]
For every multivariate Schur basis element
$s_\theta=s^{\tau_1,\ldots,\tau_k}_\phi$, including any allowed
multiplicity label, the operator of Definition 3.5 satisfies
\begin{equation}\label{eq:result}
 \G^m_\kappa s_\theta
 =-\left(\sum_{i=1}^m\phi_i(\phi_i+N+1-2i)\right)s_\theta
   +\sum_{\eta<\theta}a_{\theta\eta}s_\eta.
\end{equation}
Consequently it is strongly triangular for every choice of the
multivariate Schur basis within the specified multiplicity spaces.
\end{theorem}

\section*{An identity for the entire differential operator}
Complexify the real vector space of Hermitian matrix tuples and write
$x_{j,ab}$ for the entries of $X_j$. Thus all $km^2$ entries are independent
polynomial coordinates for the calculation; restriction to the Hermitian
space recovers the original operator. Write
$\partial_{j,ab}=\partial/\partial x_{j,ab}$ and define
\[
 \E=\sum_{j,a,b}x_{j,ab}\partial_{j,ab},\qquad
 R_{ab}=\sum_{j,c}x_{j,ca}\partial_{j,cb},\qquad
 L_{ab}=\sum_{j,c}x_{j,ac}\partial_{j,bc}.
\]
The $R_{ab}$ and $L_{ab}$ are the column and row actions of
$\mathfrak{gl}_m$, respectively. They commute with one another, and each
family satisfies the standard matrix-unit commutation relations. Set
\[
 \C_R=\sum_{a,b}R_{ab}R_{ba},\qquad
 \C_L=\sum_{a,b}L_{ab}L_{ba}.
\]
The product rule, including differentiation of the coefficient in the
inner vector field, gives
\begin{align}
 \C_R&=m\E+\sum_{j,\ell,a,b,c,d}
 x_{j,ad}x_{\ell,cb}\partial_{j,ab}\partial_{\ell,cd},\label{eq:CR}\\
 \C_L&=m\E+\sum_{j,\ell,a,b,c,d}
 x_{j,cb}x_{\ell,ad}\partial_{j,ab}\partial_{\ell,cd}.\label{eq:CL}
\end{align}
For instance, in $R_{ab}R_{ba}$, differentiation of $x_{\ell,db}$
contributes $\sum_{j,c}x_{j,ca}\partial_{j,ca}$, and summing over
$a,b$ gives $m\E$.

Define the operators that decrease degree by one,
\begin{equation}\label{eq:K}
 K_j=\sum_{a,b,c,d}\delta_{ad}x_{j,cb}
          \partial_{j,ab}\partial_{j,cd}
       +\left(\kappa_j+\frac m2\right)\sum_a\partial_{j,aa}.
\end{equation}
The source operator consists of the within-matrix second-order terms
$(\delta_{ad}-x_{j,ad})x_{j,cb}\partial_{j,ab}\partial_{j,cd}$,
the mixed terms
\[
 -\frac12\sum_{j\ne\ell,a,b,c,d}
 (x_{j,ad}x_{\ell,cb}+x_{j,cb}x_{\ell,ad})
 \partial_{j,ab}\partial_{\ell,cd},
\]
and the drift
$\sum_{j,a,b}((\kappa_j+m/2)\delta_{ab}-Nx_{j,ab})\partial_{j,ab}$.
The two quadratic coefficients coincide when $j=\ell$.
Combining them with \eqref{eq:CR}--\eqref{eq:CL} therefore proves the
exact identity, on the full polynomial ring,
\begin{equation}\label{eq:operator}
 \boxed{\quad
 \G^m_\kappa=-\frac12(\C_R+\C_L)-(N-m)\E+\sum_{j=1}^k K_j.
 \quad}
\end{equation}
Both Casimirs are retained here: they need not agree on noninvariant
polynomials. They have the same scalar action on the invariant
components used below.

\section*{The Casimir scalar, including multiplicities}
For either family of matrix units $T_{ab}$, the element
$\sum_{a,b}T_{ab}T_{ba}$ is central; this follows by substituting
$[T_{ab},T_{cd}]=\delta_{bc}T_{ad}-\delta_{ad}T_{cb}$ and cancelling
terms in its commutator with $T_{cd}$.
On the irreducible polynomial representation of highest weight
$\phi=(\phi_1,\ldots,\phi_m)$ it acts by
\begin{equation}\label{eq:casvalue}
 c_m(\phi)=\sum_{i=1}^m\phi_i(\phi_i+m+1-2i).
\end{equation}
To verify the normalization, take a highest-weight vector $v$, so that
$T_{ab}v=0$ for $a<b$ and $T_{aa}v=\phi_av$. Then
\[
 \sum_{a,b}T_{ab}T_{ba}v
 =\left(\sum_a\phi_a^2+\sum_{a<b}(\phi_a-\phi_b)\right)v
 =c_m(\phi)v.
\]
Centrality extends this equality to the whole irreducible representation.

Recall the Cauchy and tensor-product decompositions defining the source's
multivariate Schur spaces (pp.~7 and 16--17). For fixed individual
partitions $\boldsymbol\tau=(\tau_1,\ldots,\tau_k)$, each matrix contributes
a row and a column Schur factor of type $\tau_j$. Decomposing the
tensor product of row factors and the tensor product of column factors
under their respective diagonal $GL_m$ actions gives types $\phi$ and
$\psi$. Invariance under simultaneous unitary conjugation selects equal
types $\phi=\psi$. This is precisely the construction of the spaces
$S^{\boldsymbol\tau}_\phi$ in the source.

If $\phi$ has tensor-product multiplicity $d$, the corresponding invariant
space can have dimension $d^2$. Nevertheless \emph{each} copy of the
row representation has Casimir scalar $c_m(\phi)$, and the same is true
for each column copy. A Casimir belongs to the acting enveloping algebra
and acts as this scalar times the identity on the entire isotypic
component, including its multiplicity factor. Thus
\begin{equation}\label{eq:copies}
 \C_R s=c_m(\phi)s,\qquad
 \C_L s=c_m(\phi)s
 \quad\text{for every }s\in S^{\boldsymbol\tau}_\phi.
\end{equation}
This asserts more than preservation of an isotypic component by an
arbitrary intertwiner; no mixing of equivalent copies is possible.
Also $\E s=|\phi|s$, because
$|\phi|=\sum_j|\tau_j|$. The degree-preserving part of
\eqref{eq:operator} therefore acts on this entire space as
\[
 -c_m(\phi)-(N-m)|\phi|
 =-\sum_{i=1}^m\phi_i(\phi_i+N+1-2i).
\]

\section*{Every remaining term is strictly lower}
It remains to check the exact source ordering, not merely total degree.
The operator $K_j$ acts only on the entries of $X_j$, decreases their
degree by one, and is invariant under simultaneous unitary conjugation.
The last assertion follows directly from the contracted indices in
\eqref{eq:K}, or by conjugating its matrix coordinates by a unitary matrix.

Let $W_\tau$ denote the single-matrix Cauchy summand with individual
partition $\tau$. A basis element $s_\theta$ belongs to
$\bigotimes_{\ell=1}^kW_{\tau_\ell}$. If $|\tau_j|=0$, then
$K_js_\theta=0$. Otherwise
\[
 K_js_\theta\in
 \left(\bigotimes_{\ell\ne j}W_{\tau_\ell}\right)
 \otimes\mathcal P_{|\tau_j|-1}(\operatorname{Mat}_m),
\]
with factors placed in their original order. Decomposing the last
polynomial space into Cauchy summands shows that every output index
$\eta=(\rho_1,\ldots,\rho_k;\psi)$ satisfies
\[
 \rho_\ell=\tau_\ell\ (\ell\ne j),\qquad
 |\rho_j|=|\tau_j|-1.
\]
Because the source uses \emph{graded} dominance for each partition,
$\rho_j<\tau_j$, regardless of its shape. Since all the other individual
partitions agree, the source's reverse lexicographic combination gives
$\eta<\theta$, regardless of $\psi$ or its multiplicity label.
The invariant output expands in the multivariate Schur basis, so this
accounts for every term in $K_js_\theta$.

Combining this observation with \eqref{eq:operator} and
\eqref{eq:copies} proves \eqref{eq:result}. The diagonal coefficient
depends only on $\phi$, and in particular is identical on equivalent
indices. This proves strong triangularity in Definition 2.12 and
completes the proof of Theorem \ref{thm}.

\paragraph{Scope and verification.}
The argument is algebraic and establishes all matrix sizes, all numbers
of variables, all partition degrees, and all multiplicities. In fact,
the operator identity and triangularity hold for arbitrary complex
parameters; the theorem states the parameter range of the question.
The companion exact checker tests the operator identity on a basis of
quadratic polynomials in a nontrivial matrix case and verifies a
four-dimensional repeated-type invariant space. These tests check
conventions; the proof for arbitrary dimensions is above.
The review was performed by the authoring agent, without independent
peer review. This theorem is separate from the product-eigenfunction
assertion of Conjecture 3.12.

\section*{Four follow-up tasks}
\paragraph{Task 2.}
Compute the lower coefficients $a_{\theta\eta}$ explicitly in a fixed
multivariate Schur basis, starting with the first tensor-product
multiplicity greater than one, and express the general coefficients
through branching maps.
\paragraph{Task 3.}
Construct a commuting family of symmetric differential operators that
separates the partition and multiplicity labels left indistinguishable
by the single quadratic Casimir.
\paragraph{Task 4.}
Classify coincidences between the values
$\sum_i\phi_i(\phi_i+N+1-2i)$ for admissible $N$ and determine how they
affect canonical choices of orthogonal eigenpolynomials.
\paragraph{Task 5.}
Use tensor product multiplicities to compute a generating function for
the multiplicities in the Jacobi spectrum. Derive corresponding
heat-trace asymptotics in the geometric parameter range.

\begin{thebibliography}{9}
\bibitem{K} T. Kuijper, \emph{Jacobi polynomials on the simplex of Hermitian
matrices and zonal spherical functions on partial flag manifolds},
\href{https://arxiv.org/abs/2607.25760v1}{arXiv:2607.25760v1} (2026),
Definitions 2.5, 2.12, and 3.5, and Conjecture 3.13.
\end{thebibliography}
\end{document}
